% U6 WS1 -- signs, energy diagrams, heat transfer. Block 1.
\documentclass{apchem-worksheet}

\apcunit{6}
\apcunittitle{Thermochemistry}
\apcdoctitle{Worksheet 1 \textbullet\ Signs, Diagrams, and Heat Transfer}
\apcsubtitle{CED 6.1--6.3 \textbullet\ the thermometer measures the surroundings, not the system}

\begin{document}
\wsheader[Signs are written from the \textbf{system's} point of view.
Exothermic: heat leaves the system, $q < 0$, surroundings warm.
Endothermic: heat enters the system, $q > 0$, surroundings cool.]

\problem[]{\ced{6.1} Classify each and give the sign of
$q_{\text{system}}$:}
\begin{parts}
  \item Combustion of propane
        \answerblank[5.4cm]{exothermic, $q < 0$}
  \item An instant cold pack activating
        \answerblank[5.4cm]{endothermic, $q > 0$}
  \item Water freezing
        \answerblank[5.4cm]{exothermic, $q < 0$}
  \item Ice melting
        \answerblank[5.4cm]{endothermic, $q > 0$}
  \item Photosynthesis
        \answerblank[5.4cm]{endothermic, $q > 0$}
\end{parts}

\problem[]{\ced{6.1} A reaction is run in solution and the thermometer
reading rises from \SI{22.0}{\celsius} to \SI{28.4}{\celsius}.}
\begin{parts}
  \item Sign of $q_{\text{water}}$: \answerblank[2.4cm]{positive}
  \item Sign of $\Delta H_{\text{rxn}}$: \answerblank[2.4cm]{negative}
  \item Explain why the two signs are opposite even though only one event
        occurred. \answerlines[3]{They describe the same energy from
        opposite sides of the system boundary. The reaction released
        energy, so from the system's point of view $q$ is negative; the
        water absorbed that same energy, so from its point of view $q$ is
        positive.}
  \item A student writes ``the temperature went up, so the reaction is
        endothermic.'' Correct them in one sentence.
        \answerlines[2]{The thermometer was in the surroundings, not the
        system --- a rising temperature means the surroundings gained
        energy, so the reaction \emph{released} it and is exothermic.}
\end{parts}

\problem[]{\ced{6.2} Energy diagrams.}
\begin{parts}
  \item Products drawn below reactants means the reaction is:
        \answerblank[3cm]{exothermic}
  \item $\Delta H$ on a diagram is measured between:
        \answerblank[5.4cm]{reactants and products}
  \item Does the height of the barrier between them affect $\Delta H$?
        \answerlines[2]{No. $\Delta H$ depends only on the relative
        energies of reactants and products; the barrier affects the rate,
        not the enthalpy change.}
  \item Sketch, in words, the diagram for an endothermic reaction with a
        large activation barrier. \answerlines[2]{Reactants low on the
        left, a tall peak in the middle, and products at a level above the
        reactants but below the peak.}
\end{parts}

\problem[]{\ced{6.3} A hot metal block is dropped into cool water.}
\begin{parts}
  \item In which direction does energy flow?
        \answerblank[3.4cm]{hot to cold}
  \item Explain the transfer at the particle level.
        \answerlines[3]{Fast-moving atoms at the metal's surface collide
        with slower water molecules and transfer kinetic energy to them.
        The metal's particles slow on average and the water's speed up,
        which is what the two temperature changes register.}
  \item When does the transfer stop, and what is happening then?
        \answerlines[3]{Net transfer stops at thermal equilibrium, when both
        reach the same temperature. Collisions and energy exchange
        continue, but now at equal rates in both directions, so there is no
        further net change.}
  \item Why is ``the block has a lot of heat'' a poor way to put it?
        \answerlines[2]{Heat is energy \emph{in transit} because of a
        temperature difference, not a quantity a body stores. What the
        block contains is internal energy.}
\end{parts}

\problem[]{\ced{6.1} An ammonium nitrate cold pack drops from
\SI{21.0}{\celsius} to \SI{4.0}{\celsius}.}
\begin{parts}
  \item Sign of $\Delta H$ for the dissolving:
        \answerblank[2.4cm]{positive}
  \item Where did the energy come from?
        \answerlines[2]{From the water and the surroundings --- the
        dissolving process absorbed it, which is exactly why the water
        cooled.}
  \item Does an endothermic process being favourable contradict anything
        you know? Answer briefly. \answerlines[2]{No. Enthalpy is only one
        of the two factors that decide whether a process occurs; entropy is
        the other, and here the increase in disorder on dissolving is large
        enough to carry it. (Unit 9 makes this precise.)}
\end{parts}

\problem[]{\ced{6.3} Two identical blocks, one at \SI{80}{\celsius} and one
at \SI{20}{\celsius}, are placed in contact and insulated from
everything else. Predict the final temperature and justify.
\answerlines[3]{About \SI{50}{\celsius}. The blocks have the same mass and
specific heat, so equal magnitudes of energy transfer produce equal
temperature changes --- one falls by as much as the other rises, meeting in
the middle.}}

\begin{spiralstrip}{Unit 5 \textbullet\ spiral}
  \item Orders come from: \answerblank[2.6cm]{experiment}
  \item First-order half-life: \answerblank[2cm]{$0.693/k$}
  \item A catalyst changes $K$ how? \answerblank[2.6cm]{not at all}
  \item Rate-determining step is the one with:
        \answerblank[4.0cm]{the highest barrier}
  \item Molecularity gives order only for:
        \answerblank[3.8cm]{an elementary step}
\end{spiralstrip}

\end{document}
